\documentclass[a4paper, serif, 10pt]{moderncv}

%% moderncv
% color and style
\moderncvcolor{black}
\moderncvstyle{banking}

%% page layout/margins
\usepackage[top=2.5cm, bottom=2.5cm, left=1.5cm, right=1.5cm]{geometry}

\nopagenumbers{}

%% Babel config for Italian language
% ref:
% - https://latex3.github.io/babel/guides/locale-italian.html
\usepackage[italian]{babel}

%% fontspec
\usepackage{fontspec}

\IfFontExistsTF{Linux Libertine}{
  \setmainfont[Ligatures={TeX, Common}, Numbers=OldStyle]{Linux Libertine}
}{}
\IfFontExistsTF{Linux Libertine O}{
  \setmainfont[Ligatures={TeX, Common}, Numbers=OldStyle]{Linux Libertine O}
}{}

%% CTeX
% CJK support, used to show CJK characters in the resume
%
% - heading=false: disable ctex heading
% - punct=kaiming: use kaiming punctuations styles for CJK
% - scheme=plain: use plain scheme, do not override \normalsize font size
% - space=auto: space settings for CJK characters
%
% CJK fonts are resolved per locale language, e.g. Noto Serif/Sans CJK JP
% for Japanese resumes, with fallback to the next region if not installed
%
% ref:
% - http://ctan.mirrorcatalogs.com/language/chinese/ctex/ctex.pdf
\usepackage[UTF8, heading=false, punct=kaiming, scheme=plain, space=auto]{ctex}

\IfFontExistsTF{Noto Serif CJK SC}{
  \setCJKmainfont{Noto Serif CJK SC}
}{}
\IfFontExistsTF{Noto Sans CJK SC}{
  \setCJKsansfont{Noto Sans CJK SC}
}{}

%% line spacing
\usepackage{setspace}
\setstretch{1.125}

%% URL styling - use normal text instead of monospace
\urlstyle{same}

%% Patch \httplink and \httpslink to handle full URLs
% The original moderncv macros always prepend http:// or https:// to URLs.
% This causes issues when the URL already has a protocol (e.g., https://example.com)
% resulting in malformed URLs like https://https://example.com.
% This patch checks if the URL already contains "://" and uses it directly if so.
% Note: We use \str_set:Nx (x-expansion) to expand macros like \@homepage first.
\ExplSyntaxOn
\str_new:N \g_yamlresume_url_str

\RenewDocumentCommand{\httpslink}{O{}m}{
  \str_set:Nx \g_yamlresume_url_str {#2}
  \str_if_in:NnTF \g_yamlresume_url_str {://}
    { \url{#2} }
    { \url{https://#2} }
}

\RenewDocumentCommand{\httplink}{O{}m}{
  \str_set:Nx \g_yamlresume_url_str {#2}
  \str_if_in:NnTF \g_yamlresume_url_str {://}
    { \url{#2} }
    { \url{http://#2} }
}
\ExplSyntaxOff

\name{Marco Bianchi}{}
\title{Analista Dati Senior specializzato in Business Intelligence e Data Visualization}
\phone[mobile]{+39 333 123 4567}
\email{marco.bianchi@email.it}
\homepage{https://www.marcobianchi.it}

\address{Via Torino 15, Milano, Lombardia, Italia, 20121}{}{}

\extrainfo{{\small \faLinkedin}\ \href{https://www.linkedin.com/in/marcobianchi}{@marcobianchi} {} {} {} • {} {} {} 
{\small \faGithub}\ \href{https://github.com/marcobianchi}{@marcobianchi}}

\begin{document}

\maketitle

\section{Informazioni personali}

\cvline{}{\begin{itemize}
\item Analista dati con oltre 6 anni di esperienza nell'analisi di grandi dataset, modellazione statistica e creazione di dashboard interattive.
\item Esperto in Python, SQL e strumenti di BI come Power BI e Tableau, con solida conoscenza dei database relazionali e cloud (BigQuery, Snowflake).
\item Capacità dimostrata di tradurre requisiti aziendali in soluzioni data-driven, collaborando con team multidisciplinari.
\item Forte orientamento ai risultati, attenzione ai dettagli e passione per la comunicazione efficace degli insight.
\end{itemize}}
\section{Formazione}

\cventry{set 2016–lug 2018}
        {Laurea magistrale, Statistica e Data Science, Punteggio: 110/110 e lode}
        {Politecnico di Milano}
        {\url{https://www.polimi.it}}
        {}
        {\begin{itemize}
\item Tesi sperimentale su modelli predittivi per la customer churn analysis.
\item Progetto di gruppo per la creazione di una piattaforma di analisi dati in tempo reale.
\end{itemize}
\textbf{Corsi}: Analisi Statistica Multivariata, Machine Learning, Big Data Analytics, Database e Sistemi Informativi, Data Visualization, Serie Temporali}

\cventry{set 2013–lug 2016}
        {Laurea triennale, Economia e Statistica, Punteggio: 108/110}
        {Università degli Studi di Bologna}
        {\url{https://www.unibo.it}}
        {}
        {\begin{itemize}
\item Fondamenti di statistica e analisi economica.
\item Introduzione alla programmazione in R.
\end{itemize}
\textbf{Corsi}: Statistica Descrittiva, Probabilità, Econometria, Analisi dei Dati}
\section{Esperienza professionale}

\cventry{gen 2021–Presente}
        {Data Analyst Senior}
        {DataMind S.r.l.}
        {\url{https://www.datamind.it}}
        {}
        {\begin{itemize}
\item Responsabile dell'analisi dati per il reparto marketing, con focus su segmentazione clienti e ottimizzazione campagne.
\item Sviluppo di dashboard interattive in Power BI per il monitoraggio KPI aziendali, utilizzate da oltre 50 utenti.
\item Automazione di reportistica settimanale tramite Python e SQL, riducendo il tempo di produzione del 60\%.
\item Collaborazione con data engineer per la definizione di pipeline ETL e data quality.
\end{itemize}
\textbf{Parole chiave}: Power BI, Python, SQL, ETL, Marketing Analytics}

\cventry{set 2018–dic 2020}
        {Data Analyst}
        {Banca Nazionale del Lavoro}
        {\url{https://www.bnl.it}}
        {}
        {\begin{itemize}
\item Analisi di portafoglio crediti e monitoraggio rischi tramite modelli statistici in R e SQL.
\item Creazione di report mensili per la direzione, con visualizzazione di metriche di performance e anomalie.
\item Supporto alla migrazione dei dati verso un data warehouse cloud (BigQuery).
\item Partecipazione a progetti di data governance e definizione di KPI.
\end{itemize}
\textbf{Parole chiave}: R, SQL, Risk Analysis, BigQuery, Reporting}
\section{Lingue}

\cvline{Italiano}{Madrelingua o bilingue \hfill \textbf{Parole chiave}: Madrelingua}
\cvline{Inglese}{Competenza professionale completa \hfill \textbf{Parole chiave}: IELTS 8.0, Business English}
\cvline{Francese}{Competenza professionale limitata \hfill \textbf{Parole chiave}: DELF B1}
\section{Competenze}

\cvline{Analisi Dati e Statistica}{Esperto \hfill \textbf{Parole chiave}: Python, R, SQL, Statistica inferenziale, A/B testing}
\cvline{Business Intelligence}{Avanzato \hfill \textbf{Parole chiave}: Power BI, Tableau, Looker, Excel avanzato}
\cvline{Database e Cloud}{Avanzato \hfill \textbf{Parole chiave}: PostgreSQL, MySQL, BigQuery, Snowflake}
\cvline{Machine Learning}{Intermedio \hfill \textbf{Parole chiave}: scikit-learn, TensorFlow, Modelli predittivi}
\cvline{Data Engineering}{Intermedio \hfill \textbf{Parole chiave}: ETL, Apache Airflow, dbt}
\section{Premi}

\cventry{nov 2022}
        {DataMind S.r.l.}
        {Premio Innovazione Dati}
        {}
        {}
        {Riconoscimento per lo sviluppo di una soluzione di analytics predittiva che ha aumentato il tasso di conversione del 15\%.}
\section{Certificazioni}

\cventry{mar 2021}
        {Google}
        {Google Data Analytics Professional Certificate}
        {\url{https://www.coursera.org/professional-certificates/google-data-analytics}}
        {}
        {}

\cventry{giu 2022}
        {Microsoft}
        {Microsoft Certified: Data Analyst Associate}
        {\url{https://learn.microsoft.com/en-us/certifications/data-analyst-associate/}}
        {}
        {}
\section{Progetti}

\cventry{gen 2022–mar 2022}
        {Dashboard interattiva per l'analisi delle vendite retail.}
        {Dashboard Vendite Interattiva}
        {\url{https://github.com/marcobianchi/sales-dashboard}}
        {}
        {\begin{itemize}
\item Sviluppata con Power BI e dataset pubblico di vendite.
\item Include filtri dinamici, KPI e analisi trend.
\item Pubblicata su Power BI Service per la condivisione.
\end{itemize}
\textbf{Parole chiave}: Power BI, Data Visualization, Retail}

\cventry{set 2021–dic 2021}
        {Modello predittivo per la previsione del churn.}
        {Analisi Churn Clienti}
        {\url{https://github.com/marcobianchi/churn-analysis}}
        {}
        {\begin{itemize}
\item Utilizzo di Python e scikit-learn per la modellazione.
\item Feature engineering e valutazione con cross-validation.
\item Comunicazione dei risultati tramite report e presentazione.
\end{itemize}
\textbf{Parole chiave}: Python, Machine Learning, Churn}
\section{Interessi}

\cvline{Tecnologia}{Intelligenza Artificiale, Open Data, Cloud Computing}
\cvline{Sport}{Corsa, Nuoto, Ciclismo}
\section{Volontariato}

\cventry{gen 2022–Presente}
        {Mentor}
        {Data Science Milano}
        {\url{https://www.datasciencemilano.it}}
        {}
        {\begin{itemize}
\item Mentor per aspiranti data analyst, con sessioni mensili su Python e SQL.
\item Organizzazione di workshop su data visualization e carriera nel settore dati.
\item Supporto alla community nella risoluzione di problemi pratici.
\end{itemize}}
    
\end{document}